\documentclass[12pt, a4paper]{article}

\usepackage{ctex}
\usepackage{float, booktabs, parskip}
\usepackage{amsmath, amsthm, amssymb, mathrsfs, bm}
\usepackage{geometry}
\usepackage{hyperref}

\geometry{a4paper, left=2cm, right=2cm, top=2.5cm, bottom=2.5cm}

\setlength{\parindent}{2em}
\setlength{\parskip}{0.5em}
\allowdisplaybreaks[4]



\title{{\Large{\textbf{人工智能的发展历程及其在地球系统科学中的应用潜力}}}}
\author{刘滨瑞 \\ 地球系统科学系 \quad \href{mailto:lbr21@mails.tsinghua.edu.cn}{lbr21@mails.tsinghua.edu.cn}}
\date{\today}

\begin{document}

\maketitle



本报告内容基于徐冰老师于10月16日在《地球系统科学前沿》课程中的讲座《The Foundation of Artificial Intelligence and Neural Network》，旨在深入讨论人工智能的发展历程及其在地球系统科学中的应用。

报告分为三个部分。
第一章将结合讲座内容和我的个人经历，回顾人工智能从符号主义到大模型时代的范式演进。
第二章将分析当前大模型的局限性，探讨人工智能在地球系统科学中的应用潜力。
第三章将抛出一些问题，供参考。



\section{人工智能发展历程回顾}

在课程中，徐冰老师为我们详细讲述了人工智能的发展历程。
和大部分人所预想的不同，人工智能的发展并非一帆风顺，而经历了一系列剧烈的研究范式的转变。
每一次转变，该领域研究的核心问题、研究方法乃至哲学基础都发生了深刻的变化。
时至今日，人工智能早已在各行各业得到了广泛的应用，甚至已经永远地改变了人们的生活方式。
我认为，如果我们想要真正理解当下，理解这样一个“人工智能的时代”，我们就必须要理解人工智能的历史。

\subsection{符号主义、联结主义与人工智能的寒冬期}

人工智能研究的起点，实际上是对人类理性思维进行机械模拟的尝试。
1950年，图灵 (Turing) 提出了那个至今仍在引发争论的测试: \textbf{一台机器，若能在对话中使人类无法分辨其与真人的差别，我们是否应当承认它拥有智能？}
这一问题，亦被称为“图灵测试” (Turing Test)，为整个领域的研究定下了基调: 对何为“智能”的思辨。

1956年，约翰·麦卡锡 (John McCarthy) 在达特茅斯会议上正式提出“人工智能”这一术语。
随后，纽厄尔 (Allen Newell) 与西蒙 (Herbert Simon) 提出了“逻辑理论家” (Logic Theorist) 的概念，试图用形式化规则证明数学定理。
他们取得了一些成功，“逻辑理论家” 证明了罗素-怀特海《数学原理》中的一些定理。

在这一时期，研究者的核心信念是\textbf{智能即符号处理}。
他们相信，只要将世界的知识编码为符号，将推理过程表达为逻辑规则，机器就能展现智能。
这一理念，现在被称为\textbf{符号主义 (Symbolism)}。

1958年，罗森布拉特 (Frank Rosenblatt) 提出了\textbf{感知机 (Perceptron)}。
这标志着人工智能的第一次范式转变: 从符号主义转向\textbf{联结主义 (Connectionism)}。
感知机不再依赖于僵硬的符号，而是试图模仿生物神经元的连接机制。
它通过调整连接权重来学习分类任务，展现出一种“从数据中学习”的能力。
在这一时期，研究者们普遍认为智能可能并非源于精心设计的规则，而是涌现自大量简单单元的相互作用。

然而，1969年，马文·明斯基 (Marvin Minsky) 与西摩尔·派普特 (Seymour Papert) 发表了《Perceptrons》一书。
他们在数学上证明了: \textbf{单层感知机无法解决简单的异或 (XOR) 问题}。
这一理论上的致命打击，在当时的研究者们看来直接宣判了联结主义的死刑。
联结主义的“失败”，进而导致整个人工智能领域进入了长达十余年的寒冬期。
资金枯竭，研究者转向其他领域。
当然，现在的我们已经充分认识到了联结主义的价值——深度神经网络，乃是当今人工智能的基石。
但是在当时，没有人相信联结主义能够带来真正的智能。

\subsection{统计机器学习}

1980年代，\textbf{专家系统 (Expert System)}作为符号主义最后的辉煌短暂登场。
专家系统是符号主义的集大成者，它试图将专家的知识和经验编码为具体的规则，可以理解成巨量的if-then规则的集合。
MYCIN、DENDRAL等系统通过编码专家知识，在特定领域展现出了一定的实用价值。
但专家系统的局限很快暴露: 知识获取的瓶颈、过于复杂的系统设计、规则的脆弱性、几乎无法处理不确定性。
研究者们意识到，世界的复杂性难以用人为规则完全描述。

人工智能的研究方向逐渐转向了\textbf{统计机器学习}，这是第二次范式转变。
从哲学的角度，\textbf{人类承认了自身对世界的无知，我们不再试图直接去描述世界，而是选择数据自己说话}。
支持向量机 (Support Vector Machine, SVM)、决策树 (Decision Tree)、随机森林 (Random Forest)、AdaBoost等算法被提出，然后逐渐成为主流。
统计机器学习认为，\textbf{“智能”只能来源于对数据的学习 (Learning)}。
学习被重新定义为一个\textbf{优化问题 (Optimization Problem)}: 在假设空间中寻找一个函数，使其在训练数据上的经验风险最小，同时在未见数据上具有良好的泛化能力。

1984年，莱斯利·瓦利安特 (Leslie Valiant) 提出了\textbf{PAC学习框架 (Probably Approximately Correct Learning Framework)}。
PAC框架提出了一个更深刻的问题: \textbf{在计算资源和数据有限的情况下，什么问题是“可学习”的？}
通过概率分析，PAC框架将“泛化能力 (Generalization Ability)”从一个经验性的概念，转变为一个可以用数学严格证明的属性。
随着统计学习理论的进一步发展，VC维 (Vapnik-Chervonenkis Dimension) 和Rademacher复杂度 (Rademacher Complexity) 等概念被相继引入，研究者们拥有了更强大的工具来分析学习算法在有噪声数据下的泛化能力。

PAC框架具有深刻的意义，它告诉研究者们: 模型的表达能力越强，为保证其可靠的泛化能力，所需的训练样本就越多。
虽然深度神经网络的巨大成功一度让理论解释显得滞后，但统计学习理论的思想内核——通过有限的样本去逼近一个未知的真实规律，并对这一过程的可靠性给出数学保证——依然是研究者们理解和改进现代人工智能方法的根本出发点。
这为大数据时代和大模型的发展提供了理论基础，并在事实上做出了准确的预言。

\subsection{深度学习的复兴: 表示学习}

2012年，\textbf{AlexNet}，在ImageNet竞赛中取得了惊人的成绩，图像分类的正确率达到了一个前所未有的高度。
AlexNet的成功，标志着联结主义的复兴，深度神经网络开始成为人工智能的主流模型。
我认为，这一次复兴源于三个关键突破: 计算能力的飞跃 (GPU的应用)、大规模数据集的出现 (ImageNet)、以及训练技巧的改进 (ReLU、Dropout、Batch Normalization)。
其中，最重要的三位研究者是: 杰弗里·辛顿 (Geoffrey Hinton)、杨立昆 (Yann LeCun)和约书亚·本吉奥 (Yoshua Bengio)。

在这一时期，人工智能的研究范式发生了第三次转变，其核心是从\textbf{特征工程} (Feature Engineering) 转向了\textbf{表示学习} (Representation Learning)。
在统计机器学习的时代，算法的成功在很大程度上依赖于研究者手工设计的特征——例如，从原始图像中提取SIFT或HOG等描述子，或者从文本中构造TF-IDF向量。
这个过程不仅耗时耗力，更依赖于人类的先验知识和创造力，它直接构成了模型性能的瓶颈。
深度神经网络改变了游戏规则。
它不再需要人类去定义什么是好的特征，而是通过其多层级结构，直接从原始数据中自动学习有意义的表示。
例如，\textbf{卷积神经网络 (Convolutional Neural Network, CNN)} 在处理图像时，其底层网络能够自动学习到边缘和纹理等基础特征，中层网络将其组合成物体的局部部件，而高层网络则能识别出完整的物体。
对于序列数据，\textbf{循环神经网络 (Recurrent Neural Network, RNN)} 及其变体 (如LSTM和GRU) 能够捕捉时间维度上的依赖关系。
随后出现的\textbf{注意力机制 (Attention Mechanism)}，更是赋予了模型动态聚焦于输入信息中最关键部分的能力，极大地提升了处理长序列任务的性能。
这种“端到端” (End-to-end) 的学习方式，将特征提取和模型训练融为一体，极大地释放了机器学习的潜力。

这一突破开启了一个被研究者们戏称为“炼丹”的时代。
既然网络能够自动学习表示，那么研究者们只需要关心一个问题: \textbf{如何设计出更深、更强大的网络结构，以学习到更优质的表示？}
学术界的重心因此聚焦于对网络结构的精巧设计与工程优化上。
其中最重要的工作是何恺明等人提出的\textbf{残差网络 (ResNet)}。
ResNet通过引入残差连接，巧妙地解决了网络过深时出现的梯度消失和模型退化问题。
与此同时，研究者们也在不断攻克各种技术难题，从寻找更有效的优化算法 (如Adam、RMSprop)，到对抗过拟合的各种正则化技巧。

研究者们试图通过结构上的不断创新和技术上的精益求精，将模型的性能推向极限。
他们相信，更精巧的设计是通往更强智能的钥匙。

\subsection{大模型时代}

2017年，谷歌发表了论文《Attention Is All You Need》，提出了\textbf{Transformer}架构。
Transformer抛弃了循环结构，完全基于注意力机制构建。自注意力 (Self-Attention) 允许序列中的每个位置直接与其他所有位置交互，捕捉长程依赖。
多头注意力 (Multi-Head Attention) 让模型从不同角度关注输入。
位置编码 (Positional Encoding) 弥补了架构本身对顺序信息的无知。
当我第一次学习到这个架构时，我被它的简洁、高效、可并行化所震撼。

某种程度上，Transformer终结了深度神经网络的结构设计之争。
研究者们很快发现，无论是自然语言处理、计算机视觉还是语音识别，基于Transformer的简单堆叠和微调，其性能往往能超越那些为特定任务精心设计的复杂网络。
现在，Transformer已经成为了最主流的模型架构，它是一个统一、强大且可扩展性极强的框架。
从此，人工智能研究的重心不再是“发明下一个新结构”，而是“如何更好地利用和扩展Transformer”。

现在看来，Transformer的最具革命性意义的并不是其结构设计，而是在其中体现出的\textbf{尺度定律 (Scaling Law)}。
2018年，OpenAI发布了基于Transformer的第一个生成式预训练模型GPT。
随后，2019年的GPT-2展现出令人震惊的文本生成能力。
2020年，拥有1750亿参数的GPT-3则横空出世。
研究者们在实践中惊奇地发现，只要不断扩大模型规模 (参数量)、数据规模 (训练语料) 和计算规模 (GPU时长)，模型的能力就会持续增长，甚至出现训练时未曾明确优化的涌现能力 (Emergent Abilities): 小模型完全无法完成的任务，大模型突然就会了。
Kaplan等人在2020年的论文中从理论上证明了这一点: \textbf{模型性能与参数量、数据量、计算量之间存在着可预测的幂律关系}。
这意味着，只要资源充足，模型性能的提升几乎是一个必然的结果。
在Scaling Law的指引下，我们进入了\textbf{大模型时代}。
在这样一个时代，一个更简单、更粗暴的逻辑成为了主导: 极大的参数量，极大的数据量，极大的计算量，将带来具有强大智能的模型。

这一“大力出奇迹”的思路，在实践中取得了辉煌的成功。
我仍然记得第一次与GPT-3对话时的震撼: 它能流畅地回答问题、撰写文章、编写代码、进行逻辑推理。
它的知识广博、语言自然，有时展现出令人惊叹的“理解力”。
那一刻，我意识到我们正站在一个技术的临界点上。
随后的发展更是惊人，GPT-5、Gemini 2.5 Pro、Deepseek-R1等后续模型则不断刷新着人们对机器智能上限的认知。
这场由Scaling Law驱动的技术进步，其影响早已超出了学术界的范畴，引发了关于创造力、人类智能的广泛讨论。
甚至，许多人开始担心自己的工作是否会被取代。

\section{人工智能在地球系统科学中的应用}

许多研究者相信，借助尺度定律，我们正稳步走向\textbf{通用人工智能 (Artificial General Intelligence, AGI)}。
这是一种充满诱惑力的叙事，它向我们许诺了一个技术爆炸的奇点，一个机器智能超越人类的未来。
然而，一种更审慎的观点认为，仅靠“更大”的模型，并不能弥合当前范式与真正智能之间的鸿沟。
这也是我的看法。

\subsection{大模型的局限性}

Scaling Law虽然强大，但它并未改变模型能力的本质。
我认为，当前的大模型，本质上是一个卓越的\textbf{统计引擎}，而非一个真正的\textbf{推理引擎}。
其能力来源于对人类知识库的统计模式匹配，而非对物理现实的因果推理。

第一，大模型的“智能”源于对“相关性”的极致捕捉，而非对“因果性”的深刻理解。
模型在人类知识库的汪洋中学习了海量的模式与关联，从而能够生成统计上最可能的、流畅连贯的输出。
但它并不是一个真正的\textbf{世界模型 (World Model)}，不存在对物理现实、逻辑规则、因果链条的内在表征。
例如，它知道“乌云密布”和“下雨”在文本中高度相关，但它并不理解大气物理学中水汽凝结、达到饱和后形成降水的因果过程。
因此，它无法进行可靠的反事实推理——它不能真正回答“假如大气中没有凝结核，会发生什么？”这类科学问题，而只能在数据库中寻找人类是否讨论过类似的话题。

第二，大模型的知识是虚拟的，相对的，不稳固的。
模型中的每一个词，例如“海洋”，其“意义”仅仅来自于与其他词 (如“蓝色”、“波浪”、“咸水”)的统计关系。
它从未“体验”过物理世界。它能写出关于大海的优美诗篇，但它对“湿润”或“寒冷”没有任何内在概念。
这种知识与物理现实的脱节，对于日常对话或许无伤大雅，但对于科学研究却是致命的。

第三，它缺乏系统性的泛化能力。
人类智能的核心之一是组合性 (Compositionality): 我们理解了“红色”和“方块”的概念，就能立刻识别出一个此前从未见过的“红色的方块”。
一些研究指出，大模型在一定程度上具有这种能力，但还远远不足。
它擅长在训练数据的分布范围内进行插值，但当遇到需要将已知概念进行全新组合的、结构化的新问题时，它们的表现会急剧下降。
但是，科学发现恰恰要求将已知的原理应用于全新的、未曾探索过的情境中。
因此，依赖于统计模式匹配的“智能”，难以胜任科研任务。

这些局限并非简单的工程问题，而是某种固有的理论瓶颈。
当前范式所能达到的高度，终究有其极限。
我们需要一次新的理论突破——如何让机器从学习语言，转向理解语言所描述的那个“现实世界”。

\subsection{一个疑问: 大模型浪潮的可持续性}

2024年夏天，我与地学系的\textbf{韩轶轮博士后}进行了一次深入的交流。我们探讨的核心问题是: \textbf{这一轮由大模型驱动的人工智能浪潮，究竟还能持续多久？}
五年，还是更短？基于当前的技术路径和研究范式，我们近期还能期待颠覆性的突破吗？

我们认为，目前，在大模型时代，人工智能的发展路径至少会在三个方面受到严峻的挑战。

\begin{itemize}

    \item \textbf{数据}

    大模型对高质量人类数据的要求是无止境的，然而互联网上公开可及的、高质量的文本与代码数据——这座滋养了GPT系列的“富矿”——正被开采殆尽。
    更糟糕的是，随着AI生成内容 (AIGC) 的增长，未来的训练数据中，不可避免地会混入由其他模型生成的、质量参差不齐的“合成数据”。
    这可能导致“模型近亲繁殖”的退化效应，即模型开始学习自己或同类的错误与偏见，最终陷入能力停滞甚至衰退的怪圈。

    \item \textbf{经济}

    Scaling Law同时也意味着收益递减的残酷现实。
    从GPT-3到GPT-4，训练成本跃升了一个数量级，但其能力的提升，虽然显著，却远非同等比例。
    当训练下一个前沿模型的成本飙升至数十甚至上百亿美元时，仅有少数科技资本巨头才有可能参与其中了，而学术界和开源社区将逐渐被排除在外。

    \item \textbf{能源}

    训练一个旗舰级的大模型，其全生命周期的碳足迹，可能相当于数千个家庭一年的碳排放量。
    在一个日益关注气候变化和可持续发展的当下，用指数级增长的能源消耗来换取模型智能的线性增长的尝试注定是不可持续的。

\end{itemize}

一个令人警醒的结论是: 仅仅依靠“更大”的尺度，Scaling Law可能很快就要走到尽头。
在下一个理论范式出现之前，纯粹的人工智能算法研究，很有可能陷入一种“创造力的内卷”状态。
这迫使我开始思考一个不同的问题: 与其在一条可能即将停滞的赛道上等待未来，我们是否能将手中强大的工具，应用于更有潜力、更有意义的方向？

\subsection{AI4Earth: 人工智能在地球系统科学中的应用}

地球系统科学，作为关乎人类文明未来的重要领域，其内在的极端复杂性对研究者们提出了巨大的挑战。
一方面，近年来随着观测技术和数据同化技术的快速发展，地球系统科学研究也已经进入了大数据时代。
卫星遥感、地面台站、海洋浮标等观测手段每日产生数以TB计的海量、多模态数据，其中蕴藏着许多未被发掘的规律。
然而，这些数据往往充满噪声与不确定性，其巨大的体量和高维度特征也让传统统计分析方法捉襟见肘。
另一方面，地球系统是一个受物理定律严格约束的系统。
从流体力学的Navier-Stokes方程到热力学定律，能量、质量、动量守恒等基本原理，为系统的演化提供了强有力的“先验知识”。

考虑到地球系统科学\textbf{“数据丰富”}与\textbf{“规律明确”}并存的特点，我认为，人工智能在这一领域具有独特而深刻的潜力，它不仅能够充分利用海量观测数据，更能够在物理定律的严格约束下实现真正有意义的科学突破。

首先，在\textbf{高分辨率预测与模拟}方面，人工智能已经展现出了超越传统数值模式的能力。
地球系统中存在着跨越多个时空尺度的复杂相互作用: 从分钟级的对流发展到年代际的海洋环流变化，从公里级的地形影响到行星尺度的大气环流。
传统的数值天气预报模式虽然基于严格的物理方程，但受限于计算资源和网格分辨率，往往难以捕捉这些跨尺度的非线性耦合过程。
特别是对于极端天气事件——如突发性暴雨、龙卷风、热带气旋的快速增强——这些高影响、小尺度、非线性的现象，恰恰是传统模式最薄弱的环节。
深度学习模型，尤其是\textbf{时空序列模型}和\textbf{图神经网络}，能够直接从高分辨率观测数据中学习这些复杂的时空依赖关系和非线性演化模式。
它们不需要显式求解偏微分方程，而是通过端到端的学习捕捉数据中隐含的动力学规律，从而在计算效率和预测精度上实现双重突破。

其次，对于模式中的\textbf{物理过程参数化}模块，人工智能提供了一种全新的思路。
气候模型虽然包含了大尺度动力学的物理方程，但对于那些发生在网格尺度以下的物理过程——如云微物理、湍流混合、对流触发、边界层过程——必须依赖参数化方案来近似表达。
这些参数化方案往往基于简化的理论假设或经验公式，包含大量不确定的参数，成为气候模拟中最大的不确定性来源。
神经网络的强大非线性拟合能力，使其能够从高分辨率的“真实”数值模拟 (如云解析模拟) 中学习这些次网格物理过程的复杂行为，构建出比传统参数化方案更准确、更稳定的替代方案。
更重要的是，我们可以使用\textbf{物理信息神经网络 (Physics-Informed Neural Networks, PINNs)}，将物理约束嵌入神经网络的架构或损失函数中，确保人工智能学习到的参数化方案在物理上是自洽且可靠的。
我认为，这很有可能成为气候模型发展的新突破点。

最后，在\textbf{科学知识发现}这一层次上，人工智能展现出了作为科学研究工具而非仅仅是工程工具的价值。
地球系统中充满了尚未被充分理解的现象: 厄尔尼诺与全球气候的遥相关、气候系统的临界点和突变信号、极端事件的前兆特征、不同地球圈层之间的相互作用机制。
这些复杂关系往往隐藏在高维、非线性、多尺度的数据中，超出了人类专家的直接感知能力。
机器学习，特别是因果发现算法、注意力机制和可解释AI方法，能够从海量多源数据中自动挖掘出这些隐藏的气候遥相关模式，识别系统演化的因果链条，甚至发现过去从未被注意到的物理机制。
我想，人工智能可以与领域专家的物理洞察紧密结合，成为科学家的“智能助手”，帮助提出新的科学问题、指明新的研究方向。

我回忆起了软件学院\textbf{龙明盛老师}团队于2023年发表在《Nature》上的一篇研究成果——\textbf{《Skilful Nowcasting of Extreme Precipitation with NowcastNet》}。
极端降水的临近预报 (Nowcasting)，即对未来0-2小时内降水的预测，是气象学中最具挑战性的问题之一。
传统数值天气预报模式由于冷启动问题和计算时间限制，在这一时间尺度上往往力不从心。
龙老师团队设计的NowcastNet通过深度学习直接从雷达观测数据中学习降水系统的演化模式，不仅在预测精度上显著超越传统方法，更重要的是将物理约束深度融入神经网络架构——通过设计特殊的损失函数确保预测结果满足质量守恒等物理定律，使人工智能的结构本身就反映了大气物理的内在规律。
这正是\textbf{AI4Earth}的精髓所在: 并不是简单地将人工智能算法套用在地球系统科学之中，而是需要二者的深度耦合。

我很喜欢龙老师的这篇论文，这项工作对我的影响也非常大。
它让我清晰地看到了人工智能在地球系统科学中的应用潜力。



\section{若干问题探讨}

关于人工智能，以及人工智能在地球系统科学中的应用，有一些问题我一直在思考。
因此，在报告的最后一部分，我想再分享和讨论一下这些问题。

\subsection{从统计关联到因果推理}

当前人工智能在基于海量数据的统计关联性的学习方面已经取得了显著的进展。
然而，科学研究的本质，追求的并非统计上的“可能性”，而是逻辑上的“必然性”。
因此，\textbf{我们能否让人工智能从一个卓越的统计模式匹配引擎，进化为一个可靠的物理因果推理引擎？}

这要求我们超越对尺度定律的依赖，去探索神经网络与符号系统、数据驱动与知识引导相结合的混合架构。
更重要的是，我们需要建立一套针对科学应用的可信性理论与验证体系。
传统数值模式是“白盒”，其内在逻辑透明，不确定性可追溯。
而神经网络是“黑箱”，一个微小的系统性偏差就可能在长期气候模拟中被无限放大。
这在科学研究中有时是不可接受的，我们需要让人工智能的决策和预测过程变得透明，或者至少可以被验证。

\subsection{内生性的物理信息神经网络}

当前，将物理知识融入神经网络的主流方法 (如PINNs) 仍是在损失函数上施加“软约束”。
这实际上不过是一个“补丁”，强迫人工智能模型的输出符合某些物理规律。
据我所知，这样做的效果并不是很理想。
为了真正跨越“数据驱动”与“物理规律”之间的鸿沟，我们需要回答一个问题: \textbf{我们能否设计出内生性尊重物理规律的新型神经网络架构？}
例如，设计出其结构天然满足特定守恒定律 (如质量守恒、能量守恒) 的网络层，或者将对称性、不变性等基本物理原则，变为网络自身的归纳偏置。
这需要更先进的数学工具的帮助，例如辛几何、李代数等。

\subsection{大模型时代的教育}

作为一个在大模型时代求学的大学生，我亲身经历了人工智能所带来的知识获取方式的变革。
我很庆幸，我的基础知识体系是在没有大语言模型辅助的情况下，通过艰苦的思考、阅读和调试建立起来的。
这个过程是痛苦的，但是，它塑造了我的批判性思维，也培养了我独立解决问题的能力。
现在，学生们可以轻易地从大语言模型那里获得问题的答案。
这大大提高了学习的效率，但是，如果学生们从一开始就依赖于大语言模型，是否会导致他们永远错过一些特定的能力培养的机会？

我能够高效地使用人工智能，将大语言模型作为我的“助手”，恰恰是因为我拥有独立完成任务的能力和足够的知识积累。
这使得我可以判断模型给出的答案是否正确，是否存在幻觉；然后，我可以提出更精准、深刻的问题。
但是，假如一个学生从一开始就依赖于人工智能呢？他可能会成为大语言模型的高效使用者，可是这还远远不够。

我认为，知识的传授是容易的，但是能力的培养是困难的。
人工智能可以传授知识，但无法替代那个艰难却必要的、将知识内化为能力的过程。
这意味着，在人工智能的时代，我们迫切地需要一种全新的教育方式。

\end{document}

